\documentclass{article}
\usepackage{float}
% Language setting
\usepackage[german]{babel}
\usepackage[table]{xcolor}
% Set page size and margins
% Replace `letterpaper' with `a4paper' for UK/EU standard size
\usepackage{multicol}
\usepackage[letterpaper,top=2cm,bottom=2cm,left=3cm,right=3cm,marginparwidth=1.75cm]{geometry}
\usepackage[dvipsnames]{xcolor}
% Useful packages
\usepackage{amsmath} 
\usepackage{graphicx}
\usepackage[colorlinks=true, allcolors=blue]{hyperref}
\usepackage{subcaption}
\usepackage{graphicx}
\title{Test 2 - Fragen auf Kolloquien} 
\author{Lil'Jana}
\usepackage{enumitem,amssymb}
\newlist{todolist}{itemize}{2}
\setlist[todolist]{label=$\square$}
\usepackage{pifont}




\begin{document}
\maketitle
\subsection*{Die DNA, die im Gentechnik-Labor bearbeitet wird, wird meist zunächst zerkleinert durch ...}
\begin{itemize}     \item  \end{itemize}
\subsection*{Was sind die wesentlichen Unterschiede zwischen DNA und RNA?}
\begin{itemize}
    \item 
    \item 
\end{itemize}
\subsection*{Welche Eigenschaften sollte ein Vektor mindestens haben?}
\begin{itemize}
    \item 
    \item 
    \item    
\end{itemize}
\subsection*{Wofür steht die Abkürzung 'PCR'?}
\begin{itemize}
    \item 
\end{itemize}

\subsection*{Wie lauten die vier fehlenden Bezeichnungen?}
\begin{figure}[H] 
    \centering
    \includegraphics[width=10cm]{bilder/nukleotif-fehlend.png}
    \caption{}
    \label{fig:my_label}
\end{figure}
\subsection*{Die Funktion von Restriktionsenzymen in Bakterien ist ....}
\begin{itemize}     \item  \end{itemize}
\subsection*{Homologe Chromosomen unterscheiden sich häufig in den Merkmals-Informationen, die sie tragen. Sie ...}
\begin{itemize}     \item  \end{itemize}
\subsection*{Die tRNA, die die Aminosäure Valin translatiert, hat das Anticodon CAG.
Wie lautet das dazugehörige mRNA-Triplett?}
\begin{itemize}     \item  \end{itemize}
\subsection*{In welchem Meiosestadium lagern sich die homologen Chromosomen zusammen, und findet "crossover" statt?}
\begin{itemize}     \item  \end{itemize}

\subsection*{Der Mensch hat 46 Chromosomen. Diese Anzahl (menschlicher)
Chromosomen findet man in ...}
\begin{itemize}     \item  \end{itemize}

\subsection*{Ein Karyogramm wird hergestellt, in dem man ...}
\begin{itemize}     \item  \end{itemize}

\subsection*{Was ist ein Diplont?}
\begin{itemize}     \item  \end{itemize}

\subsection*{Wenn Sie Crossover bzw. den Austausch von genetischem Material in
einer Tetrade zwischen Chromatiden berücksichtigen, was wird in der
Meiose I voneinander getrennt?}
\begin{itemize}
    \item 
\end{itemize}
\subsection*{Eine Pflanze mit dem Genotyp AABbcc ist ...}
\begin{itemize}
    \item 
\end{itemize}
\subsection*{Manchmal kann ein Gen die Auspägung eines anderen Gens verhindern
oder überdecken. Diese genetische Wechselwirkung bezeichnet man als}
\begin{itemize}
    \item 
\end{itemize}


\subsection*{Pleiotropie ist eine Genwirkung, bei der ...}
\begin{itemize}
    \item 
\end{itemize}

\subsection*{Wie heißt die Genwirkung, bei der viele Gene auf ein Merkmal wirken?}
\begin{itemize}
    \item 
\end{itemize}

\subsection*{Formel zur Berechnung der Rekumbinationsfrequenz}

\subsection*{Die an einem RNA-Template (Matrize) synthetisierte DNA nennt man}
\begin{itemize}
    \item 
\end{itemize}




\subsection*{Bei der Erstellung von genetischen (Gen-) Karten werden aus Daten zu
Rekombinationsfrequenzen zwischen Genen genetische Distanzen
errechnet und die Abstände in cM angegeben. Was ist die Ursache für
dafür, dass Rekombinationsfrequenz-Werte (in $\%$) nicht einfach additiv in
Centi-Morgan umgerechnet werden können?}
\begin{itemize}
    \item \bigskip
\end{itemize}

\subsection*{Welche Bedingungen muss eine Population bzw. die Allele in der
Population erfüllen, damit das Hardy-Weinberg-Gesetz für die
Allelfrequenzen in Populationen gültig ist?}
\begin{itemize}
    \item 
    \item 
    \item 
    \item 
\end{itemize}

\subsection*{Was für eine Mutation (welche das/die fragliche(n) Gen(e) betrifft) wäre
wahrscheinlich am wenigsten schädlich für das zu bildende Protein?}
\begin{itemize}
    \item 
\end{itemize}

\subsection*{Den Austausch von Chromosomenabschnitten zwischen nicht-homologen
Chromosomen bezeichnet man als ...}
\begin{itemize}
    \item 
\end{itemize}

\subsection*{Was ist Inversion?}
\begin{itemize}
    \item 
\end{itemize}
\subsection*{Was ist Tranduktion?}
\begin{itemize}
    \item 
\end{itemize}

\subsection*{Eine Person, die den Genotyp XXY aufweist, hat ....}
\begin{itemize}
    \item 
\end{itemize}



\subsection*{Wie viele (protein-kodierende) Gene befinden sich in etwa auf dem Y-Chromosom des Menschen?}
\begin{itemize}
    \item 
\end{itemize}

\subsection*{Die snRNPs genannten Partikel sind ...}
\begin{itemize}
    \item 
\end{itemize}

\subsection*{Definieren Sie 'Spleißosom'}
\bigskip\bigskip \bigskip
\subsection*{Normalerweise werden bei eukaryotischen Genen ...}
\begin{todolist}% %
\item die Exons nicht transkribiert.  
\item die Introns nicht transkribiert.  
\item die Introns zwar in RNA transkribiert, aber diese RNA verlässt den
Zellkern nicht.  
\item die Exons zwar in RNA transkribiert, aber diese RNA verlässt den
Zellkern nicht.  
\item Exons und Introns in RNA transkribiert, und diese RNA verlässt den
Zellkern.  
\end{todolist}

\subsection*{Beispiele für Totipotenz sind:}
\begin{itemize}
    \item 
\end{itemize}
\subsection*{Definieren Sie 'Totipotenz'}
\begin{itemize}     \item  \end{itemize}
\bigskip\bigskip \bigskip

\subsection*{Eine einzelne embryonale Stammzelle kann sich nicht zu einem
kompletten Embryo entwickeln, weil ...}
\bigskip\bigskip 


\subsection*{Wodurch werden die Körperachsen in der frühen Entwicklung von D.
melanogaster festgelegt?}
\bigskip\bigskip 

\subsection*{Apoptose ist ...}
\bigskip\bigskip 

\subsection*{Bei Eukaryoten bezeichet man bestimmte Strukturen im Kern, die aus ca.
200 bp DNA und Protein bestehen, als ....}
\bigskip
\subsection*{Was ist für den Unterschied in der Genomgrösse zwischen Arabidopsis
thaliana und Zea mays die wesentliche Ursache?}
\bigskip\bigskip 

\subsection*{Wie kann man sich die Ausbreitung von Retrotransposons in Genomen
erklären?}
\bigskip\bigskip\bigskip\bigskip


\subsection*{Die Embryonalentwicklung von Pflanzen und Tieren ...}
\bigskip\bigskip

\subsection*{Wie schreibt man DNA-Sequenzen?}
\bigskip\bigskip

\subsection*{Wohin wandert DNA im elektrischen Feld?}
\begin{itemize}     \item  \end{itemize}
\bigskip\bigskip

\subsection*{Welche strukturellen Eigenschaften haben DNA und Polyphosphat gemeinsam? }
\begin{itemize}
    \item 
    \item 
    \item 
    \item 
\end{itemize}

\subsection*{Wie färbt man DNA?}
\begin{itemize}
    \item 
    \item 
\end{itemize}


\subsection*{Wie unterscheidet sich DNA von RNA?}
\begin{itemize}
    \item 
    \begin{itemize}
        \item DNA  
        \item RNA  
    \end{itemize}
    \item 
    \begin{itemize}
        \item DNA  
        \item RNA  
    \end{itemize}
    \item 
    \begin{itemize}
        \item DNA  
        \item RNA 
    \end{itemize}
    \item 
    \begin{itemize}
        \item DNA: 
        \item RNA: 
    \end{itemize}
    \item 
    \begin{itemize}
        \item DNA 
        \item RNA 
    \end{itemize}
    \item 
    \begin{itemize}
        \item DNA: 
        \item RNA: 
    \end{itemize}
\end{itemize}

\subsection*{Warum ist DNA stabiler als RNA?}

\begin{itemize}
    \item 
    \item 
    \item 
    \item 
\end{itemize}

\subsection*{Welcher Phosphatrest wird in die DNA eingebaut?}
\begin{itemize}     \item  \end{itemize}

\subsection*{Welcher Phosphatrest wird bei Kinase-Reaktionen transferiert?}
\begin{itemize}     \item  \end{itemize}
 
\subsection*{Was kann man über die Kinetik der DNA-Synthese sagen?}
\begin{itemize}
    \item 
    \item 
    \item 
    \item 
    \item 
    \item
\end{itemize} 

\subsection*{DNA-Replikation}
\subsubsection*{1. Ist wie gerichtet?} 
\begin{itemize}     \item  \end{itemize}

\subsubsection*{2. Verhält sich wie bezüglich der Ausgangs-DNA und Tochter-DNA?} 
\begin{itemize}     \item  \end{itemize}

\subsubsection*{3. Verhält sich wie bezüglich des Ablaufs der Replikation an den beiden DNA-Strängen?} 
Es gibt einen führenden und einen verzögerten Strang.
\begin{table}[H]
    \centering
    \begin{tabular}{|p{7cm}|p{7cm}|} \hline
     führend & verzögert \\ \hline \hline
    \begin{itemize}     \item  \end{itemize}\begin{itemize}     \item  \end{itemize}\begin{itemize}     \item  \end{itemize} &   \begin{itemize}     \item  \end{itemize}\begin{itemize}     \item  \end{itemize}\begin{itemize}     \item  \end{itemize} \\\hline
\end{tabular}
\end{table} 



\subsection*{Geben Sie die Werte der Parameter der DNA-
Superhelikalitat an? Wie erfolgen die Umwandlungen?}
\subsubsection*{Formel von LK, T und W:}

\begin{figure}[H] 
    \centering
    \includegraphics[width=7cm]{bilder/linkerausfüllen.png}
    \label{fig:my_label}
    \caption{}
\end{figure}
\subsection*{Aufgabe: In einer Probe seien DNA und RNA enthalten. Sie sollen nur die RNA zerstören. Wie gehen Sie vor?}
\begin{itemize}     \item  \end{itemize}


\subsection*{Welche(n) Parameter der Superhelikalität der DNA können nur Topoisomerasen verändern?}
\begin{todolist}
\item  Linking Number (LK)
\item Verdrilling Twist (T) 
\item Anzahl der Überkreuzungen Writhe (W)
\end{todolist}
\subsection*{Wie unterscheidet man superhelikale von 'entspannter' DNA?}
\begin{table}[H]
    \centering
    \begin{tabular}{|c||p{5cm}|p{5cm}|} \hline 
         & Superhelikal & Entspannt  \\ \hline\hline
        Verdrillung & \begin{itemize}     \item  \end{itemize} & \\ \hline
        Topologische Eigenschaften&\begin{itemize}     \item  \end{itemize}& \\ \hline
        Funktionelle Bedeutung &\begin{itemize}     \item  \end{itemize}& \\ \hline
    \end{tabular}
    \label{tab:my_label}
\end{table} 
\subsection*{Lückentext}
Bei der Replikation der DNA von E. coli arbeitet die DNA-Polymerase III. sehr genau. Sie synthetisiert nur jedes \textbf{.........................}te Basenpaar falsch, weil sie selbst über eine Reparaturfunktion verfügt, die \textbf{.........................}  ist. \newline
Replizierte DNA weist jedoch noch weniger Fehler auf: nur jedes \textbf{.........................}  te Basenpaar falsch, weil es zusätzliche Reparaturmechanismen gibt, z.B. die \textbf{.........................} .

\subsection*{Wie können Mutationen induziert werden?}
\subsubsection*{1. Physikalisch, z.B.} 
    
    \begin{itemize}
    \item 
    \item 
    \item
\end{itemize}
\subsubsection*{2. Chemisch, z.B.} 
\begin{itemize}
    \item 
    \item 
    \item 
\end{itemize}


\subsection*{Die Promotor Regionen...}
\centerline{AAAAAATCTTGACACCCAGCATCAGTAGACATAATATGACCACA\textbf{T}GGACCAAA} 

\subsection*{Die RNA-Polymerase ist wesentlich für die Transkription in E. coli. Welchem Enzym der DNA-Replikation (Primase oder DNA-Polymerase III) ähnelt sie mehr?}
Primase und RNA-Polymerase. 
\begin{table}[h] 
    \centering
    \begin{tabular}{|p{3cm}||p{4cm}|p{4cm}|p{4cm}|}\hline
         & Primase &  DNA-Polymerase III & RNA-Polymerase\\ \hline \hline
         Aktivität /was wird synthetisiert & & &  RNA \\\hline
         Vorläufer/Substrate &  & & NTPs \textit{(Ribonukleosidtriphosphate)} \\\hline
         Matrize & & &  ssDNA \\\hline
         Starter oder Primer & & &  Kein Primer benötigt \\\hline
    \end{tabular}
    \label{tab:my_label}
\end{table}

\subsection*{Wie viele verschiedene Aminosäuren gibt es in
Proteinen?}
\begin{itemize}     \item  \end{itemize}
\begin{itemize}     \item  \end{itemize}


\subsection*{Ordnen Sie die
Translationsfaktoren EF-G,EF-TU, IF1, IF2, IF3, RF und RRF zeitlich in den Prozess der Translation ein und weisen Ihnen eine Funktion dabei zu. 
}
\begin{enumerate}
    \item  ........................................
    \begin{enumerate}
        \item ........................................
        \item ........................................
        \item ........................................
    \end{enumerate}
    \item ........................................
    \begin{enumerate}
        \item ........................................
        \item ........................................
    \end{enumerate}
    \item ........................................
    \begin{enumerate}
        \item ........................................
        \item ........................................
    \end{enumerate}
\end{enumerate}
\subsection*{Was ist  die Shine-Dalgarno-Sequenz?
Wo liegt sie? Was ist  ihre Funktion?}
\begin{itemize}     \item  \end{itemize}\begin{itemize}     \item  \end{itemize}

\subsection*{Wie sind Ribosomen aus Archaea und Eukaryoten zusammengesetzt?} 
\begin{table}[H]
    \centering
    \begin{tabular}{ |p{2cm}|| p{4cm}| p{4cm}| p{4cm}|}
    \hline
     & Archaea & Bacteria & Eukaryoten \\ \hline \hline
    Typ &   &  & \\ \hline
    zusammen gesetzt aus..  &   &  & \\ \hline
    rRNA &  &   &  & \\ \hline
    ribosomale Proteine  &   &  & \\ \hline
    Hemmstoffe &   &  & \\ \hline
\end{tabular}
\end{table}

\subsection*{Was ist das katalytisch aktive Prinzip des Ribosoms?}
\begin{itemize}     \item  \end{itemize} \bigskip

\subsection*{Aufgabe: Ordnen Sie die Vorgänge bei der homologen Rekombination. (nicht alle existieren)}
Branch migration Wandern der Holliday-Struktur Eindringen eines Einzelstrangs
Eindringen eines Doppelstrangs Auf losen der Holliday-Struktur Doppelstrang-Bruch
Ausbildung von 3'-Überhängen Ausbildung von 5'-Überhängen
\begin{enumerate}
    \item ........................................
    \item ........................................
    \item........................................
    \item ........................................
    \item ........................................
\end{enumerate}


\subsection*{Bei der Rekombination zweier homologer DNA-Regionen mit den nach einander liegenden Genen 1, 2, 3 und 4 bildet sich eine Holliday-Struktur aus, die über die Gene 2 und 3 hinweg wandert. Zeichnen Sie die Schritte der Rekombination und lösen Sie die Holliday-Struktur auf zwei Arten auf.} 
\newpage

\subsection*{Bei welchen Genen kann Genkonversion auftreten?}
\begin{itemize}     \item  \end{itemize}
\subsection*{Als Oswald Avery zu einem Transformationsansatz mit einem Gemisch aus durch Hitze getöteten 5-Zellen und lebenden R-Zellen von Streptococcus pneumoniae folgende Substanz zusetzte unterblieb die Transformation und die mit dem Gemisch infizierte Maus überlebte.
}
\begin{itemize}     \item  \end{itemize}
\subsection*{Bei Bakterien gibt es verschiedene Möglichkeiten der DNA-Übertragung. Der direkte Zell-Zell-Kontakt ist wesentlich bei der ...}
\begin{itemize}     \item  \end{itemize}
\subsection*{Sie transformieren einen Tryptophan-, Arginin- und Leucin- auxotrophen Bakterien-Stamm trp- arg- leu- mit DNA eines Arginin- und Leucin-prototrophen, aber Tryptophan-auxotrophen Stammes trp - arg + leu+. 
Welche Rekombinante(n) kann/können entstehen?}
\begin{itemize}     \item  \end{itemize}

\subsection*{Enzyme und Proteine der DNA-Replikation}
\begin{table}[H]
    \centering 
    \begin{tabular}{|p{4cm}|p{11cm}|}
    \hline
     &  Initiation der Replikation am oriC\\\hline 
     & Bindet an einzelsträngige DNA  \\\hline
    Helikase &  \\\hline
    Primase &   \\\hline
     & Synthetisiert den neuen DNA-Strang \\\hline
     & Entfernt den RNA-Primer (5'-3'-Exonuklease) und füllt den Strang auf \\\hline
    Ligase &  \\\hline
    Topoisomerase &  \\\hline
\end{tabular}
\end{table}
\subsection*{tRNA}
\subsubsection*{Wie lang sind tRNA-Moleküle?}
\begin{itemize}     \item  \end{itemize}
\subsubsection*{Welche Struktur nehmen tRNA-Moleküle ein?}
\begin{itemize}     \item  \end{itemize}
\subsubsection*{Was kann man über die Basen in tRNA sagen?}
\begin{itemize}     \item  \end{itemize}  \bigskip

\subsubsection*{Welche Funktion hat tRNA?}
\begin{itemize}     \item  \end{itemize}
\subsection*{Beschreiben Sie das trp-Operon und seine Kontrolle}
\begin{table}[H]
    \centering
    \begin{tabular}{|c||p{5cm}|p{5cm}|} \hline
         & Repression und Aktivierung
der Transkriptionsinitiation &  Transkriptionelle
Attenuation \\ \hline \hline
        Protein & Repressor oder Aktivator &   \\ \hline 
         Transkription & Initiation wird reguliert &   \\ \hline 
        Vorkommen & Prokaryoten und Eukaryoten &  \\ \hline 
    \end{tabular}
    \label{tab:my_label}
\end{table}
\subsection*{Sie wollen ein Protein mit einem Wirtsbakterium produzieren. Dazu klonieren
Sie das Gen für dieses Protein mitsamt der für Transkription, Translation und Sekretion notwendigen Sequenzen. Welches Bakterium wählen Sie?}
\begin{itemize}     \item  \end{itemize}


\subsection*{Bei der Klonierung von Genen nutzt man gerne die a-Komplementation.
Wieviele Gene dafür sind wo kodiert?}
\begin{itemize}     \item  \end{itemize}

\subsection*{Wo und warum setzt man die Taq Polymerase ein?}
\begin{itemize}     \item  \end{itemize}
\subsection*{Beschreiben Sie Taq Polymerase bezüglich ...}
\subsubsection*{... der Substrate}
\begin{itemize}     \item  \end{itemize}
\subsubsection*{... der Matrize}
\begin{itemize}     \item  \end{itemize}
\subsubsection*{... der Aktivität}
\begin{itemize}     \item  \end{itemize}
\subsection*{Lückentext}
Bei der Expressionsklonierung können Sie experimentell der Frage nach dem ................... wissenschaftsphilosophischen Kriterium nachgehen. 
\subsection*{Welche Elemente muß die zu klonierende DNA aufweisen?}
\begin{itemize}
    \item 
    \item 
    \item 
    \item 
\end{itemize}
\subsection*{Beschreiben Sie die Schritte der Gram-Färbung!}
Beschreiben Sie die Schritte der Gram-Färbung!
\begin{enumerate}
    \item 
    \item 
    \item 
    \item 
\end{enumerate}

\subsection*{SSU rDNA-Sequenzanalysen erlauben es alle Lebewesen der Welt phylogenetisch zu betrachten. Was ist die Einheit als evolutionärer Abstandsmesser, was als evolutionäre Uhr?}
\textbf{“Evolutionäre Abstandsmesser”} \newline \begin{itemize}     \item  \end{itemize}
\textbf{“Evolutionäre Uhr”} \newline \begin{itemize}     \item  \end{itemize}
\subsection*{Der phylogenetische Stammbaum der Lebewesen basiert auf SSU rDNA-Sequenzanalysen. Beschriften sie die 3 Reiche und zeichnen Sie die ungefähren Positionen von E. coli, Mensch, Maus und Mais ein!}
\begin{figure}[H]
    \centering
    \includegraphics[width=10cm]{bilder/maisposition.png}
    \caption{}
    \label{fig:my_label}
\end{figure}
\subsection*{Was katalysiert Lysozym?}
\begin{itemize}     \item  \end{itemize}
\subsection*{Was katalysiert Transpeptidase?}
\begin{itemize}     \item  \end{itemize}
\subsection*{Was katalysiert Transglykosidase?} 
\begin{itemize}     \item  \end{itemize}
\subsection*{Wie wirkt Penicillin (und andere beta-Lactam-Antikiotika)?}
\begin{itemize}
    \item 
    \item 
    \item 
\end{itemize}
\subsection*{Skizzieren Sie die Zellwand ...}
\subsubsection*{... eines Gram-positiven Bakterium}
\bigskip\bigskip\bigskip\bigskip\bigskip
\subsubsection*{... eines Gram-negativen Bakteriums}
\bigskip\bigskip\bigskip\bigskip\bigskip

\subsection*{Was versteht man
in der mikrobiellen Ökologie unter „great plate
count anomaly“?}
\begin{itemize}     \item  \end{itemize}
\subsection*{Wie setzt sich das bakterielle Murein
zusammen? Wie ist es verknüpft?}
\subsubsection{Zusammensetzung}
\begin{itemize}
    \item 
    \item 
\end{itemize}
\subsubsection{Verknüpfung}
\begin{itemize}
    \item 
\end{itemize}
\subsection*{Resistenz gegen Penicillin}
\subsubsection*{Wie kann ich Penicillin-resistente Bakterien
selektionieren?}
\begin{itemize}     \item  \end{itemize}
\subsubsection*{Welches Enzym kann Penicillin-Resistenz
vermitteln?}
\begin{itemize}     \item  \end{itemize}
\subsection*{Gel-Elektrophorese}
\subsubsection*{Welchen Typ der Gel-Elektrophorese nutzt
man zur Analyse und Größenfraktionierung
von DNA?}
\begin{itemize}     \item  \end{itemize}
\subsubsection*{Welchen Typ der Gel-Elektrophorese nutzt
man zur Analyse und Größenfraktionierung
von Proteinen?}
\begin{itemize}     \item  \end{itemize}
\subsection*{Was ist das Prinzip der DNA-
Sequenzierung nach Sanger?}
\begin{itemize}     \item  \end{itemize} \bigskip\bigskip
\subsection*{Beschreiben Sie die
Genaktivierung durch NtrC~P}
\begin{enumerate}
    \item 
\item 
\item 
\end{enumerate}

\subsection*{Listen Sie die Aminosäuren}
\begin{enumerate}
    \item \item \item \item 
    \item \item \item \item 
    \item \item \item \item 
    \item \item \item \item 
    \item \item \item \item 
\end{enumerate}


\end{document}